\documentclass[11pt]{amsart}
\usepackage{amsmath,amssymb,amsthm}
% A small note written for Madalean's tests and demo (CC0).
\newcommand{\N}{\mathbb{N}}
\newcommand{\Z}{\mathbb{Z}}
\DeclareMathOperator{\parity}{par}
\theoremstyle{plain}
\newtheorem{theorem}{Theorem}[section]
\newtheorem{lemma}[theorem]{Lemma}
\newtheorem{proposition}[theorem]{Proposition}
\newtheorem{corollary}[theorem]{Corollary}
\newtheorem{conjecture}[theorem]{Conjecture}
\theoremstyle{definition}
\newtheorem{definition}[theorem]{Definition}
\newtheorem*{remark}{Remark}

\title{A note on parity}
\author{Ada Example}

\begin{document}
\begin{abstract}
We recall why the sum of two even integers is even, prove that the square of an even integer is even,
and state an open-looking conjecture about odd squares.
\end{abstract}
\maketitle

\section{Definitions}
Throughout, integers are elements of $\Z$ and natural numbers are elements of $\N$.
Parity questions go back to Euclid \cite{euclid}; see also the textbook of Hardy and Wright \cite{hardywright}.

\begin{definition}\label{def:even}
An integer $n$ is \emph{even} if there is an integer $k$ such that $n = 2k$, and \emph{odd} if $n = 2k+1$.
\end{definition}

\begin{lemma}[Closure]\label{lem:sum}
If $a$ and $b$ are even, then $a + b$ is even.
\end{lemma}
\begin{proof}
By Definition~\ref{def:even}, write $a = 2k$ and $b = 2l$. Then
\begin{equation}\label{eq:sum}
a + b = 2(k + l),
\end{equation}
which is even.
\end{proof}

\section{Squares}
The next result is the main point of this note.

\begin{theorem}\label{thm:square}
If $n$ is even, then $n^2$ is even.
\end{theorem}
\begin{proof}
Write $n = 2k$ as in Definition~\ref{def:even}. Then $n^2 = 2(2k^2)$, and we conclude with
Lemma~\ref{lem:sum} applied to $a = b = 2k^2$, using \eqref{eq:sum}.
\end{proof}

\begin{corollary}\label{cor:fourth}
If $n$ is even, then $n^4$ is even.
\end{corollary}
\begin{proof}
Apply Theorem~\ref{thm:square} to $n$ and then to $n^2$.
\end{proof}

\begin{remark}
The converse of Theorem~\ref{thm:square} also holds \cite{hardywright}.
\end{remark}

\begin{conjecture}\label{conj:odd}
Every odd square is congruent to $1$ modulo $8$.
\end{conjecture}

\begin{thebibliography}{9}
\bibitem{euclid} Euclid, \emph{Elements}, Book IX.
\bibitem{hardywright} G. H. Hardy and E. M. Wright, \emph{An introduction to the theory of numbers}, Oxford, 1979.
\end{thebibliography}
\end{document}
